\section{Introduction}

Geographic information frictions concentrate retail demand for locally headquartered stocks \citep{covalMoskowitz1999, ivkovicWeisbenner2005}, and that demand moves prices \citep{hongKubikStein2008}. If state-level financial sophistication shapes how retail investors trade --- slower to chase recent winners, more attentive to volatility --- then state literacy should moderate cross-sectional return anomalies that depend on retail mispricing. The natural test is the interaction between momentum and idiosyncratic volatility (IV): limits to arbitrage scale with IV \citep{shleiferVishny1997, pontiff2006}, local retail demand is sophistication-dependent \citep{calvetCampbellSodini2009, grinblattKeloharjuLinnainmaa2012}, and geographic ownership shapes anomalies more broadly \citep{bernileKumarSulaeman2015}. This paper asks whether NFCS Big-Five state literacy \citep{lusardiMitchell2014} dampens the momentum~$\times$~IV interaction for U.S. common stock, and where in the cross-section the dampening lives.

The headline finding is that the moderation concentrates in the middle of the institutional-ownership (IO) distribution. On a 519{,}090 firm-month / 6{,}081 firm sample of CRSP-listed U.S. common stock between 2009 and 2023, restricted to firms whose Compustat headquarters state is stable across the panel, stratifying the three-way regression by within-stock IO tercile produces a mid-tercile coefficient of $\hat\gamma_{\text{mid}} = -0.0321$ (state-clustered $t = -2.63$; wild-cluster bootstrap $p = 0.064$; family-wise Romano-Wolf $p = 0.077$). The low- and high-IO coefficients are roughly one-third and one-sixth of that magnitude. A continuous-margin quadratic-IO specification corroborates the U-shape with an interior trough at $\mathrm{IO}^{\ast} \approx 0.34$ ($\hat\gamma_{5} = +0.0151$, $p = 0.017$). Literacy dominates the nearest alternative state variable: a horse race against state college attainment within mid-IO returns $\hat\gamma_{\text{lit}} = -0.0350$ ($t = -2.97$) and $\hat\gamma_{\text{college}} = +0.00076$ ($t = +0.65$).

Three bounds discipline the result. First, the quadratic clears five percent but does not survive log-IO, kink-IO, or a twenty-bin diagnostic; the discrete-tercile contrast is the functional-form-free characterization, and the quadratic provides only parametric corroboration. Second, a point-in-time literacy reassignment built from EDGAR 10-K headers --- which correctly handles the 14 percent of firm-months whose then-current headquarters state differs from the static Compustat snapshot --- collapses the aggregate three-way coefficient by roughly thirty-nine fold ($-0.0117 \to -0.0003$) and the mid-IO coefficient by twenty-six fold ($-0.0321 \to -0.00121$), and flips the sign of the within-stock retail-versus-institutional difference. The directional within-tercile pattern survives on the corrected full panel; the magnitudes do not. The corrected sample refutes the prior aggregate prediction; the documented mid-IO concentration on stable-HQ replaces it as the contribution. Third, relocators are nearly twice as concentrated in the low-IO tercile as elsewhere, so the stable-HQ mid/low ratio of $2.83$ overstates the within-stock contrast relative to the corrected full-panel ratio of $1.09$. About one-third of that gap is low-IO attenuation and two-thirds is mid-IO growth.

The result is a characterization, not a microfoundation. The mid-IO concentration discriminates against any generic state-level moderator that hits all headquarters-state stocks alike, but does not select a single within-family channel. The canonical reading --- that limits to arbitrage drive the moderation through IV-scaled arbitrage capacity --- predicts that the literacy effect within mid-IO strengthens at higher IV deciles. It does not: a within-month NYSE IV-decile decomposition of the mid-IO three-way fails to reject equality across deciles (joint $\chi^{2}$ $p = 0.94$ persistent, $p = 0.18$ time-varying), and decile 10 returns $\hat\gamma \approx 0$. The data therefore rule out the canonical arbitrage-capacity microfoundation, but the IV-decile null cannot distinguish compositional heterogeneity within mid-IO from a no-channel null. I name compositional heterogeneity as the most direct family-internal description of what remains rather than as a positively-identified channel. The quadratic-IO curvature is stable across the 2015 NFCS-wave break under a Chow test, but the test is uninformative --- its contrast standard error is roughly ninety percent of $\hat\gamma_{5}$ itself --- and the trough drifts from $0.26$ to $0.40$ as the median IO share rises (Section~\ref{sec:mech-stability}).

The closest prior work is \citet{korniotisKumar2013}, who show that state-level business cycles predict returns of locally headquartered firms. This paper differs on three margins: the state variable is NFCS literacy on a three-year wave frequency rather than the quarterly business-cycle frequency; the object of study is a three-way moderation of an anomaly interaction rather than a level prediction; and the discriminating margin is a within-stock IO contrast that the state business-cycle channel does not predict. A business-cycle control block attenuates the aggregate literacy three-way by 20--23 percent, so the two state moderators are distinct but partly overlapping; the discrimination claim's scope is generic state-level confounds, not ownership-interactive ones, which remain open. \citet{korniotisKumarPage2020} use brokerage-demographic proxies for state-level sophistication, but neither NFCS literacy nor the momentum-IV interaction. \citet{hillertJacobsMuller2014} is the closest precedent for a state-level moderator of momentum, using media coverage as the primary moderator and state investor individualism as a cross-sectional moderator of the media-momentum effect, rather than literacy.
